\documentclass[12pt]{article}

\usepackage[margin = .8in]{geometry}
\usepackage{amsmath}
\usepackage{graphicx}
\usepackage{multicol, enumerate, tabularx}

\usepackage{adjustbox, soul, setspace}

\usepackage{fancyhdr}
\pagestyle{fancy}

\renewcommand{\emph}[1]{\textsf{\textbf{#1}}} %%% bold emph


\lhead{Math F113X: Numbers and Society}
\rhead{Date: \hspace{1in}}

\usepackage{tikz}
\usetikzlibrary{calc,trees,positioning,arrows,fit,shapes,through, backgrounds}
\usetikzlibrary{patterns}

\usetikzlibrary{decorations.markings}
\usetikzlibrary{arrows}

\usepackage{pgfplots}

\usepackage{longtable}
\usepackage{tabularx}

\newcommand{\ds}{\displaystyle}
\newcommand{\ans}[1][1in]{\rule{#1}{.5pt}}

\newcommand{\points}[1]{(#1 points.)}		% Trying to be lazy.

\usepackage{array}
\newcolumntype{L}[1]{>{\raggedright\let\newline\\\arraybackslash\hspace{0pt}}m{#1}}
\newcolumntype{C}[1]{>{\centering\let\newline\\\arraybackslash\hspace{0pt}}m{#1}}
\newcolumntype{R}[1]{>{\raggedleft\let\newline\\\arraybackslash\hspace{0pt}}m{#1}}
\newcommand{\red}[1]{\textcolor{red}{#1}}

\newcommand{\be}{\begin{enumerate}}
\newcommand{\ee}{\end{enumerate}}

\begin{document}
\begin{center}
{\Large  Worksheet Finance 3: Interest}
\end{center}



%\noindent \textbf{Group names:} \hrulefill \\
%-------------------------------------------------------------------------------------------------------------
%						Assignment
%-----------------------------------------------------------------------------------------------------



\begin{enumerate}
%%%S 2.1.6 probs 11 and 12
\item The purpose here is to highlight the importance of reading carefully. Subtle difference in wording can mean very different things.\\

For this problem, assume an investment has an initial deposit of \$1600 and no other deposits or withdrawals are made. \\

\textbf{Your work should include the expression or calculation you put into the calculating device } as well as the numerical answer.
	\begin{enumerate}
	\item 5\% of the initial deposit is \\
	
	\vfill
	
	\item The account increased by 5\% means the account gained \$ {\hrulefill}.\\
	
	\vfill

	
	\item The account increased by 5\% means the new balance in the account is \\
	
	\vfill

	
	
	\item The new balance is 105\% of the original balance means the new balance is \\
	
	\vfill

	
	
	\item Your final numerical answers in parts (c) and (d) should be the same though your work/calculations should be different. Answer the following by doing both calculations.\\
	
	The account increased by 20\% means the new balance is \\
	
	\vfill

	
	The account increased by 100\% means the new balance is \\
	
	\vfill

	
	\item What is a more common way of saying an account ``increased by 100\%"?\\
	
	\vfill

	
	\item Suppose you are told that each month the new balance is 104\% of the old balance, what percent of the old balance is \emph{gained} each month?\\
	
	\vfill

	
	\end{enumerate}
\newpage
%2.1.6 \# 13
\item This problem will compare simple interest and compound interest. (Formulas below.)

$$\text{simple interest:  } A=P \ast(1+\left(\frac{r}{n}\right)\ast t) \hspace{1in} \text{compound interest:  } A=P \ast\left(1+\frac{r}{n}\right)^{(n\ast t)}$$

You deposit of \$2000 into a savings account that pays 6\% interest quarterly. Set up a Simple Interest versus Compound Interest calculator. (See below.)\\

\vspace*{-2.3in}
\noindent \includegraphics[angle=-90,scale=0.6]{finance3pic.pdf}

\vspace*{-2in}

	\begin{enumerate}
	\item What are the numerical entries in cells \verb'E3' and \verb'F3' and what do they represent?

	\vfill

	\item For each account (simple and compound), calculate how much of the balance after 5 years is \emph{interest}?
	\vfill

	\item For each account (simple and compound), what is the \emph{percent increase} in the account over 5 years and how does this compare to the interest rate?

	\vfill
	
	

	\item Change the entry in \verb'D3' from a 5 to a 30 and answer questions (b) and (c) in this case.

	\vfill
	
	\item Keeping $t=30,$ change $P$ to $4000$ and determine the percent increase. 
	\vfill
	
	\item Make some observations about your previous  calculations.
	
	\vfill
	
	\vfill

\end{enumerate}

\end{enumerate}
\end{document}
\item Use a spreadsheet to calculate the following quantities:
\be
\item $471+5619$ = \ans
\item $36*987$= \ans
\item $45^{3}$= \ans (Use syntax \verb`=45^3`)
\item $12!$= \ans (Use syntax \verb`=fact(12)`)
\item Use a spreadsheet to calculate your total bill if the dinner cost \$72.50 and you want to give an 18\% tip. \ans

\ee

\singlespacing

\item Use a spreadsheet to calculate the following if you invested \$4500  for 10 years, assuming \emph{simple interest}.  (Use the process from the lecture notes.)

\onehalfspacing
\be
\item how much interest would you earn after 10 years if your interest rate was 2\%? \ans

\item How much total money did you have at the end of 10 years? \ans

\item How much interest would you earn after 10 years if your interest rate was 3\%? \ans

\item How much total money did you have at the end of 10 years? \ans

\item How much \emph{additional interest} did you earn in increasing the rate from 2\% to 3\%? \ans

\ee

\singlespacing
\item Use a spreadsheet to calculate the following if you invested \$4500  for 10 years, assuming \emph{compound interest}, compounded annually, instead. Do the calculation in the same sheet as the previous problem, so you can compare.

\onehalfspacing

\be
\item How much total interest would you have earned after 10 years if your interest rate was 2\%? \ans

\item How much total money did you have at the end of 10 years? \ans

\item How much total interest would you have earned after 10 years if your interest rate was 3\%? \ans

\item How much total money did you have at the end of 10 years? \ans

\item How much \emph{additional interest} did you earn in increasing the rate from 2\% to 3\%? \ans

\item How much more money did you make over the 10 years by using compound interest rather than simple interest, if the interest rate was 2\%? \ans \ 3\%? \ans

\ee

\newpage

\singlespacing
\item Now we will make a more clever interest calculator. 

\be[i.]
\item Make a new sheet in your spreadsheet, called \verb`Interest`.
\item In \verb`A1` type \verb`Simple Interest`
\item In row 2 make cells A through E contain the words 

\begin{tabular}{c c c c c}
\fbox{ {\tt year}} &\fbox{ {\tt principal}}&	\fbox{ {\tt interest}} &	\fbox{ {\tt total interest}} &	\fbox{ {\tt grand total}} 
\\
\end{tabular}
\item Start column {\tt A2} with 1 and then 2 and fill down until year 10.
\item In {\tt C1} type \fbox{{\tt 0.06}} (this is where we are storing our interest)
\item Type \fbox{{\tt 500}} into {\tt B3}
\item Type \fbox{{\tt =B\$3*C\$1}} into {\tt C3} 
\item In {\tt D3} type \fbox{{\tt  =sum(C\$3:C3)}} 
\item In {\tt E3} type \fbox{{\tt =B\$3+D3}} 
\item Fill down {\tt C3, D3, E3} until year 10
\ee

 To convert the long decimals to things that look like currency, you can click on the column header and then click the button that looks like \verb`$` in the toolbar.

\doublespacing

\be
\item How much total \emph{interest} had you earned by Year 10? \ans
\item How much total \emph{money} had you earned by Year 10? \ans
\item Fill down more columns if necessary. At what year will you have more than \$1200? \ans
\item At what year will you have earned more than \$500 in interest in total? \ans
\ee



\singlespacing

\item Answer the previous questions assuming you invested \$825 at 2.75\% interest using simple interest:

\doublespacing

\be
\item What two cells do you need to update to make this change?

Cell  \ans\ should change to \ans

 Cell \ans\ should change to \ans


 

\item How much total \emph{interest} had you earned by Year 10? \ans
\item How much total \emph{money} had you earned by Year 10? \ans
\item Fill down more columns if necessary. At what year will you have more than \$1200? \ans
\item At what year will you have earned more than \$500 in interest in total? \ans

\ee


\newpage

\singlespacing
\item We will next make a more clever compound interest calculator by copying the simple interest calculator and updating some cells. 

\be[(i)]
\item Select your simple interest spreadsheet cells, copy them, click on cell {\tt G1} and hit paste.
\item Change the title to \emph{Compound Interest}
\item Change your cell values so  {\tt I1}  contains \fbox{{\tt 0.06}} and {\tt H3}  contains \fbox{{\tt 500}}.
\item Type \fbox{{\tt =H3*I\$1}} into {\tt I3}. 
\item Type \fbox{{\tt =H3+I3}} into {\tt H4}.
\item Type \fbox{{\tt =H4*I\$1}} into {\tt I4}.
\item Now select cells {\tt H4, I4} and fill them down. Columns {\tt J} and {\tt K} should update automatically.
\item If you like, select columns {\tt H,I,J,K} and turn them into currency by clicking the \$ in the toolbar.
\ee

Answer the following questions about compound interest, investing \$500 compounded annually at a rate of 6\%:



\be
\item How much interest did you earn in Year 10? \ans
\item How much total money had you earned by Year 10? \ans
\item How much more money did you have in Year 10 than what you earned from simple interest? \ans
\item Fill down more columns if necessary. At what year will you have more than \$1200 in total? \ans
\item How did this answer differ from your simple interest answer? \ans
\item At what year will you have earned more than \$500 in total interest? \ans
\item How did this answer differ from your simple interest answer? \ans


\ee

\item Now answer the previous questions about compound interest assuming you invested \$825 at 2.75\% interest:
\be
\item What two cells do you need to update to make this change?

Cell  \ans\ should change to \ans

 Cell \ans\ should change to \ans


\item How much interest did you earn in Year 10? \ans
\item How much total money had you earned by Year 10? \ans
\item Fill down more columns if necessary. At what year will you have more than \$1200 in total?\ans
\item At what year will you have earned more than \$500 total in interest? \ans

\ee
 
 \ee


\end{document}

%-------------------------------------------------------------------------------------------------------------------------------------------------------------------------------------------------------------------

%%% Local Variables:
%%% mode: latex
%%% TeX-master: t
%%% End:
